\subsection{Tangent spaces, tangent linear maps}

5.5.1. Let X be a manifold and let \(a\in\mathbb{X}\) . Consider the pairs \((c,h)\) where \(c=(\mathrm{U},\varphi,\mathrm{E})\) is a chart of the manifold X at \(a\) and where \(h\) is an element of E. Two such pairs \((c,h)\) and \((c',h')\) are said to be equivalent if the derivative at \(\varphi(a)\) of the map \(\varphi'\circ\varphi^{-1}\) (which is defined on a neighborhood of \(\varphi(a)\) ) transforms \(h\) into \(h'\) . One thus obtains an equivalence relation between pairs \((c,h)\) and one calls a tangent vector at \(a\) to X a class of equivalent pairs \((c,h)\) (Sets, Chap.~II, § 6, no. 9).

The tangent vectors at \(a\) to X form a set denoted \(\mathrm{T}_{a}(\mathrm{X})\) . If \(c = (\mathrm{U}, \varphi, \mathrm{E})\) is a chart of the manifold X at \(a\) , the map \(\theta_{c}\) from E into \(\mathrm{T}_{a}(\mathrm{X})\) which associates to an element \(h\) of E the tangent vector represented by the pair \((c, h)\) is a bijection. If we transport to \(\mathrm{T}_{a}(\mathrm{X})\) by means of \(\theta_{c}\) the structure of topological K-vector space on E, the structure thus obtained does not depend on the choice of \(c\) and turns \(\mathrm{T}_{a}(\mathrm{X})\) into a Banach space, called the tangent space to X at \(a\) . The dimension (finite or \(+\infty\) ) of \(\mathrm{T}_{a}(\mathrm{X})\) is equal to the dimension of X at \(a\) .

5.5.2. Let X, Y be two varieties, \(f\) be a morphism from X to Y and \(a\) be a point of X. Consider a chart \(c = (U, \varphi, E)\) of X at \(a\) and a chart \(c' = (V, \psi, F)\) of Y at \(b\) with \(f(U) \subset V\) ; the map \(\Phi = \psi \circ f \circ \varphi^{-1}\) of \(\varphi(U)\) to \(F\) is of class \(C^{r}\) , and its derivative \(\mathbf{D}\Phi(\varphi(a))\) at the point \(\varphi(a)\)

is a continuous linear map from E into F. The continuous linear map \(\theta_{c'} \circ D\Phi(\varphi(a)) \circ \theta_{c}^{-1}\) of \(T_{a}(X)\) to \(T_{b}(Y)\) does not depend on the charts \(c\) and \(c'\) chosen; it is denoted \(T_{a}(f)\) and is called the tangent linear map to \(f\) at \(a\) . If \(g\) is a morphism from Y into a manifold Z, we have:

\[
\mathrm{T} _ {a} (g \circ f) = \mathrm{T} _ {f (a)} (g) \circ \mathrm{T} _ {a} (f).
\]

5.5.3. Let \(f:X\to Y\) be a morphism of manifolds. If \(f\) is locally constant, we have \(\mathrm{T}_{a}(f)=0\) for all \(a\) in X. Conversely, if \(\mathrm{T}_{a}(f)=0\) for all \(a\in X\) and if the field K has characteristic 0, then \(f\) is locally constant.

5.5.4. Let U be an open subset of a manifold X and let f be the canonical injection from U into X; for every point a in U, the map \(\mathrm{T}_{a}(f)\) of \(\mathrm{T}_{a}(\mathrm{U})\) to \(\mathrm{T}_{a}(\mathrm{X})\) is an isomorphism of topological vector spaces, by which these two spaces will henceforth be identified.

5.5.5. Let U be an open subset of a Banach space E; if \(\iota\) is the canonical injection of U into E, the triple \(c = (U, \iota, E)\) is a chart of U, and the atlas \(\{c\}\) defines the manifold structure of U. Given a point \(a\) of U, the map \(\theta_c\) is an isomorphism of topological vector spaces from E onto \(T_a(U)\) ; the inverse isomorphism will be denoted \(\zeta_E(a)\) .

Let \(f\) be a map of class \(\mathbf{C}^{r}\) from U into an open subset V of a Banach space F; for every point \(a\) of U, we have:

\[
\mathbf {D} f (a) \circ \zeta_ {\mathbf {E}} (a) = \zeta_ {\mathbf {F}} (f (a)) \circ \mathrm{T} _ {a} (f)
\]

where Df(a) is the derivative of f at a (1.2.1). In other words, the diagram:

\[
\begin{array}{c} \mathrm{T} _ {a} (\mathrm{U}) \xrightarrow {\mathrm{T} _ {a} (f)} \mathrm{T} _ {f (a)} (\mathrm{V}) \\ \Biggl \downarrow \zeta_ {\mathbf {E}} (a) \qquad \qquad \qquad \Biggl \downarrow \zeta_ {\mathbf {F}} (f (a)) \\ \mathrm{E} \xrightarrow {\mathrm{D} f (a)} \mathrm{F} \end{array}
\]

is commutative.

5.5.6. A cotangent vector to X at a, also called a tangent covector or simply a covector at a, is any continuous linear form on the topological vector space \(\mathrm{T}_{a}(\mathrm{X})\) ; these covectors are therefore the elements of the topological dual \(\mathrm{T}_{a}(\mathrm{X})'\) of the tangent space to X at a. This space will be endowed with the topology of uniform convergence on bounded subsets of \(\mathrm{T}_{a}(\mathrm{X})\) ; equipped with this topology, \(\mathrm{T}_{a}(\mathrm{X})'\) is a Banach space called the cotangent space to X at a. It is also denoted by \(\mathrm{T}_{a}^{\prime}(\mathrm{X})\) .

Let \(f\) be a morphism from X into a Banach space E. One calls the differential of \(f\) at \(a\) and denotes by \(d_{a}f\) or \(df_{a}\) the continuous linear map \(\zeta_{\mathbf{E}}(f(a)) \circ \mathbf{T}_{a}(f)\) of \(\mathbf{T}_{a}(\mathbf{X})\) into E. For \(t \in \mathbf{T}_{a}(\mathbf{X})\) , one sometimes denotes \(t(f)\) or t.f the value \(d_{a}f(t)\) of the linear map \(d_{a}f\) at the point t. The map \(f \mapsto d_{a}f\) is linear.

If \(E = K\) , the differential \(d_a f\) of \(f\) at \(a\) is a covector of \(X\) at \(a\) .

Let E, F, G be three Banach spaces and let \((u, v) \mapsto u \cdot v\) be a continuous bilinear map from \(F \times G\) into E. Let f (resp. g) be a morphism from X to F (resp. G). Then the map \(f \cdot g : x \mapsto f(x) \cdot g(x)\) is a morphism from X to E and we have, for every \(t \in \mathrm{T}_{a}(\mathrm{X})\) :

\[
d _ {a} (f. g) (t) = f (a). d _ {a} g (t) + d _ {a} f (t). g (a).
\]

Taking in particular G = E and F = K, the bilinear map in question being multiplication, we have:

\[
d _ {a} (f g) = f (a) d _ {a} g + g (a) d _ {a} f.
\]

5.5.7. Let X be a locally finite-dimensional manifold and let \(\xi^{1},\ldots,\xi^{m}\) be morphic functions defined in an open neighborhood U of a point \(a\) of X. The following conditions are equivalent:

\begin{enumerate}
\def\labelenumi{(\roman{enumi})}
\item
  there exists an open neighborhood V of a contained in U, such that the functions \(\xi^{i}|V\) (for \(1 \leqslant i \leqslant m\) ) form a coordinate system of X in V;
\item
  the differentials \(d_a\xi^i\) for \(1 \leqslant i \leqslant m\) form a basis of \(T_a(X)'\) ;
\item
  the map \(\xi = (\xi^1, \ldots, \xi^m)\) from U into \(K^m\) is étale at \(a\) (see no. 5.7.6).
\end{enumerate}

For the differentials \(d_{a}\xi^{1},\ldots,d_{a}\xi^{m}\) to be linearly independent, it is necessary and sufficient that there exist a neighborhood V of \(a\) , contained in U, such that the functions \(\xi^{i}|V\) are part of a coordinate system of X in V.

5.5.8. Let X be a locally finite-dimensional manifold and let \(\xi = (\xi^1, \ldots, \xi^m)\) be a coordinate system of X in an open set U. Let \(a \in \mathbf{U}\) : we denote by \((\partial_{1,a}, \ldots, \partial_{m,a})\) the basis of the tangent space \(\mathrm{T}_a(\mathrm{X})\) dual to the basis \((d_a\xi^1, \ldots, d_a\xi^m)\) of \(\mathrm{T}_a(\mathrm{X})'\) . We therefore have:

\[
\partial_ {i, a} (\xi^ {j}) = \delta_ {i, j} \quad (\text { Kronecker index }).
\]

The tangent vector \(\partial_{i,a}\) is also denoted \((\partial/\partial\xi^{i})_{a}\) .

Let \(f\) be a function of class \(\mathbf{C}^r\) on U, with values in a Banach space E. We denote by \(\partial_i f\) or \(\partial f / \partial \xi^i\) the function \(a \mapsto \partial_{i,a}(f)\) : this is a function of class \(\mathbf{C}^{r-1}\) on U with values in E (a continuous function if \(r = 1\) ). Its value at a point \(a\) of U is sometimes denoted \((\partial f / \partial \xi^i)_a\) . For every \(t \in \mathrm{T}_a(X)\) , we have:

\[
d _ {a} f (t) = \sum_ {i = 1} ^ {m} d _ {a} \xi^ {i} (t) \frac {\partial f}{\partial \xi^ {i}} (a),
\]

which can also be written:

\[
d _ {a} f = \sum_ {i = 1} ^ {m} \frac {\partial f}{\partial \xi^ {i}} (a) d _ {a} \xi^ {i}.
\]

These notations are consistent with those of 1.6.3.

5.5.9. Let X and Y be two manifolds and let f be a morphism from X to Y. We call the rank of f at a point a of X, and denote by \(rg_{a}f\) , the rank (finite or equal to +∞) of the linear map \(\overline{T}_{a}(f)\) . The map \(a \mapsto rg_{a}f\) is lower semicontinuous.

